\documentclass[10.5pt]{article}

\usepackage[utf8]{inputenc}
\usepackage[T1]{fontenc}
\usepackage[margin=0.6in]{geometry}
\usepackage{enumitem}
\usepackage{titlesec}
\usepackage{xcolor}
\usepackage{hyperref}
\hypersetup{pdfborder={0 0 0}}
\usepackage{helvet}

% Clickable link: plain underlined text, same color/font as surrounding
% text, so it reads as clickable without introducing color or boxes.
\newcommand{\linkbadge}[2]{%
  \href{#1}{\underline{#2}}%
}

\renewcommand{\familydefault}{\sfdefault}
\pagestyle{empty}

\definecolor{rulegray}{HTML}{444444}

\titleformat{\section}
  {\large\bfseries}
  {}{0em}{}
  [{\vspace{1pt}\titlerule}]
\titlespacing{\section}{0pt}{7pt}{3pt}

\newcommand{\cventry}[2]{%
  \par\noindent\textbf{#1}\hfill\makebox[1.4in][r]{#2}\par
}
\newcommand{\cvsub}[1]{%
  \par\noindent\textit{#1}\par\vspace{2pt}
}

\setlist[itemize]{leftmargin=14pt, itemsep=1pt, topsep=0pt, parsep=0pt, partopsep=0pt}

\begin{document}

\begin{center}
  {\LARGE \textbf{Nisong Monyimba}}\\[3pt]
  1515 S Extension Rd, Mesa, AZ 85210 \;|\; nmonyimb@asu.edu \;|\; 602-545-6789 \;|\;
  \linkbadge{https://github.com/NisongMonyimba/AppliedMathThesis}{github.com/NisongMonyimba/AppliedMathThesis}
\end{center}
\vspace{2pt}

\noindent Applied mathematician with a background in biomedical engineering whose research focuses on
stochastic analysis, mean-field games, stochastic control, and scientific computing. Current work develops
analytical and computational methods for common-noise mean-field games, combining stochastic differential
equations, optimal transport, numerical analysis, and reproducible research software. Seeking to pursue this
work through doctoral study in applied mathematics.

\section*{Research Interests}
\noindent Stochastic Analysis \textbullet\ Mean-Field Games \textbullet\ Numerical Analysis \textbullet\
Scientific Computing \textbullet\ Stochastic Control \textbullet\ Optimal Transport \textbullet\
Partial Differential Equations

\section*{Education}
\cventry{Biomedical Engineering, M.Sc. --- Arizona State University}{Graduated May 2022}
\cvsub{Ira A. Fulton Schools of Engineering; MasterCard Foundation Scholar; Dean's List: Fall 2020, Spring 2021}
\cventry{Biomedical Engineering, B.Sc. --- KNUST}{Graduated May 2021}
\cvsub{Kwame Nkrumah University of Science \& Technology; MasterCard Foundation Scholar; First Class Honours}

\noindent \textbf{Undergraduate quantitative coursework:} Calculus, Differential Equations, Numerical Methods,
Probability and Statistics, Signals and Systems.

\section*{Independent Mathematical Preparation}
\cventry{Graduate-Level Self-Study}{2025--Present}
\cvsub{Undertaken independently of formal coursework, in connection with the research program below}
\begin{itemize}
  \item Real Analysis, Functional Analysis, Measure-Theoretic Probability, Stochastic Calculus, Mean-Field
  Games, McKean--Vlasov SDEs, Optimal Transport, and Numerical PDEs
\end{itemize}

\section*{Research}

\cventry{Independent Research Program in Applied Mathematics}{2025--Present}
\cvsub{Mean-field games \textbullet\ stochastic control \textbullet\ Wasserstein geometry \textbullet\ optimal
transport \textbullet\ FBSDEs \textbullet\ numerical PDEs \textbullet\ Lean 4 formal verification}
\noindent Current research develops mathematical and computational methods for common-noise mean-field games,
combining stochastic analysis, numerical analysis, scientific computing, and reproducible research software.
Current outputs: a theoretical preprint, a numerical preprint, a software repository, and an independent
research monograph (below).

\vspace{2pt}
\noindent\textbf{Publications and Preprints}

\vspace{2pt}
\noindent Monyimba, N. (2026). \textit{Quantitative Contraction Estimates for Common-Noise Mean-Field Games.}
Preprint. Establishes a local contraction estimate, on an explicit time horizon, for the fixed-point map
characterising equilibria of mean-field games with common noise, extended to arbitrary horizons via a
continuation argument.\\
\linkbadge{https://github.com/NisongMonyimba/AppliedMathThesis/blob/main/papers/paper1_wellposedness/main.pdf}{PDF}
\;|\;
\linkbadge{https://github.com/NisongMonyimba/AppliedMathThesis/tree/main/papers/paper1_wellposedness}{GitHub}
\;|\; DOI: \linkbadge{https://doi.org/10.5281/zenodo.21131132}{10.5281/zenodo.21131132}

\vspace{2pt}
\noindent Monyimba, N. (2026). \textit{Numerical Approximation of Common-Noise Mean-Field Games via Particle
Methods and Finite Differences.} Preprint. Combined Euler--Maruyama, particle-method, and finite-difference
scheme for common-noise mean-field game equilibria, with every reported convergence rate the direct,
reproducible output of executed code.\\
\linkbadge{https://github.com/NisongMonyimba/AppliedMathThesis/blob/main/papers/paper2_numerics/main.pdf}{PDF}
\;|\;
\linkbadge{https://github.com/NisongMonyimba/AppliedMathThesis/tree/main/papers/paper2_numerics}{GitHub}
\;|\; DOI: \linkbadge{https://doi.org/10.5281/zenodo.21131132}{10.5281/zenodo.21131132}

\vspace{2pt}
\noindent\textbf{Independent Research Monograph}

\vspace{2pt}
\noindent\textit{Quantitative Convergence Analysis of Stochastic Maximum Principles for Mean-Field Games with
Common Noise} (2026). A book-length treatment developed independently following the 2022 M.Sc., covering
well-posedness, particle methods, finite-difference schemes, and partial Lean 4 formalisation.\\
\linkbadge{https://github.com/NisongMonyimba/AppliedMathThesis/blob/main/manuscript/main.pdf}{PDF}
\;|\;
\linkbadge{https://github.com/NisongMonyimba/AppliedMathThesis}{Full repository}

\vspace{2pt}
\noindent\textbf{Current Research Directions}
\begin{itemize}
  \item Nonlinear (non-LQ) common-noise mean-field games
  \item Numerical stochastic control and common-noise master equations
  \item Numerical methods for high-dimensional FBSDEs and PDEs
  \item Formal verification of stochastic analysis results in Lean 4
\end{itemize}

\section*{Research Software}

\cventry{AppliedMathThesis (Python, Lean 4)}{}
\cvsub{Research software implementing numerical methods for common-noise mean-field games}
\begin{itemize}
  \item SDE solvers (Euler--Maruyama, Milstein), including mean-field/common-noise variants
  \item Particle methods with conditional propagation-of-chaos verification
  \item Finite-difference Fokker--Planck solver
  \item Wasserstein-distance convergence verification and automated numerical benchmarks
  \item Reproducible computational experiments; version-controlled, continuously integrated (GitHub Actions), automatically tested (pytest)
  \item Lean 4 formalisation of core estimates (Gronwall, contraction bound, propagation of chaos)
  \item \linkbadge{https://github.com/NisongMonyimba/AppliedMathThesis}{github.com/NisongMonyimba/AppliedMathThesis}
  \;|\; DOI: \linkbadge{https://doi.org/10.5281/zenodo.21131132}{10.5281/zenodo.21131132}
\end{itemize}

\section*{Research Experience}

\cventry{Research Assistant, Weaver Lab, Arizona State University}{2021--2022}
\begin{itemize}
  \item Analyzed trophoblast proliferation dynamics by quantifying cell counts over time and fitting growth curves.
  \item Used Fiji to extract quantitative metrics from laboratory images, including colony expansion area and
  fluorescence intensity, supporting rigorous modeling of cellular growth patterns.
  \item Quantified cell proliferation rates from fluorescence microscopy data using Fiji and GraphPad Prism;
  developed statistical plots and regression analyses to characterize proliferation and viability.
  \item Developed Python-based probabilistic simulations and Monte Carlo analyses to evaluate variability in
  material properties and stochastic loading, improving predictive accuracy.
\end{itemize}

\cventry{Research Assistant, Dr. Pizziconi's Lab, Arizona State University}{2022--2023}
\begin{itemize}
  \item Designed and optimized bone fixation screws using Halpin--Tsai micromechanics models to predict composite
  behavior under mechanical load.
  \item Created and refined 3D CAD models in SolidWorks, ensuring dimensional accuracy, manufacturability, and
  integration with experimental validation workflows.
  \item Trained and mentored lab members in modeling, simulation, and experimental validation, connecting
  theoretical mathematics with engineering practice.
  \item Collaborated with interdisciplinary teams to iterate designs from simulation results; documented and
  presented research outcomes supporting publication efforts and potential clinical translation.
\end{itemize}

\cventry{Research Assistant, Medical Device Regulation (MDR) Project, KNUST}{2020--2021}
\begin{itemize}
  \item Analyzed regulatory frameworks governing medical devices across 20 African countries, identifying
  disparities in approval processes, safety standards, and market accessibility.
  \item Used publicly available regulatory data to categorize countries by regulatory maturity, from fully
  established authorities to emerging or informal frameworks.
\end{itemize}

\cventry{Med Tech R\&D Engineer, Mayo Clinic \& Arizona State University}{2022--2023}
\begin{itemize}
  \item Designed bone fixation screws using Halpin--Tsai equations and Mori--Tanaka models.
  \item Ran nonlinear FEA simulations to analyze fatigue, stress distribution, and deformation behavior.
\end{itemize}

\section*{Teaching Philosophy}
\noindent I aim to teach mathematics by combining rigorous theory with computational examples and
reproducible scientific programming.

\section*{Teaching Experience}

\cventry{Teaching Interests}{}
\cvsub{Based on completed undergraduate coursework, self-directed graduate-level study, and the teaching and
tutoring experience below}
\begin{itemize}
  \item Calculus I \& II, Linear Algebra, Differential Equations, Numerical Analysis, Probability \& Statistics,
  Engineering Mathematics, Scientific Computing (Python)
\end{itemize}

\cventry{Teaching Assistant, KNUST}{2019}
\begin{itemize}
  \item Tutored biomedical engineering undergraduates in calculus, electrical circuits, and biostatistics;
  assisted with lab supervision, grading, and practical demonstrations.
\end{itemize}

\cventry{Science \& Math Teacher, NSMQ Team, Kwanyako Senior High School}{2016--2018}
\begin{itemize}
  \item Prepared and led students for regional and national science and math competitions, emphasizing teamwork,
  critical thinking, and problem-solving.
\end{itemize}

\cventry{Math Teacher, Witness Academy}{2016--2017}
\begin{itemize}
  \item Taught core mathematics and physics; prepared junior-high students for final exams (BECE).
\end{itemize}

\section*{Master's Applied Project \& Undergraduate Capstone Project}

\cventry{Master's Applied Project: Trophoblast-on-Chip Microfluidic Device}{2022}
\begin{itemize}
  \item Investigated trophoblast growth to quantify proliferation dynamics (cell counts over time, growth curves,
  proliferation index) using Fiji to extract quantitative growth metrics from laboratory cell images.
  \item Performed statistical analysis including growth-curve comparison across conditions (two-way ANOVA),
  doubling-time comparison between treatments (unpaired t-test), and proliferation-marker quantification (Kruskal--Wallis).
  \item Developed microfluidic channel designs in SolidWorks, simulated laminar flow and shear profiles in COMSOL
  Multiphysics, and validated models against analytical equations.
\end{itemize}

\cventry{Capstone Project: Novel ECMO Cannula Design, Arizona State University}{2021}
\begin{itemize}
  \item Designed and simulated ECMO cannula flow in COMSOL Multiphysics, optimizing shear stress and velocity
  distribution; built a 3D model in SolidWorks and simulated flow in ANSYS Fluent.
  \item Quantified flow uniformity and reduced clot-formation risk by optimizing inlet geometry and flow-rate parameters.
  \item Performed finite element analyses (FEA) in ANSYS and SolidWorks to assess performance and safety margins
  under physiological load conditions.
\end{itemize}

\section*{Internship Experience}

\cventry{REU Internship, Case Western Reserve University}{2020}
\begin{itemize}
  \item Studied biological design and product development principles; designed Arduino-based air-quality sensors
  and prototyped PCB systems using EAGLE CAD.
  \item Collaborated with interdisciplinary teams to apply quantitative modeling to life-science problems.
\end{itemize}

\cventry{BME Internship, Komfo Anokye Teaching Hospital, Kumasi, Ghana}{2019}
\begin{itemize}
  \item Supported medical equipment maintenance and calibration; performed voltage and current testing.
  \item Documented device performance data and proposed engineering solutions to minimize error variance;
  presented findings to senior biomedical engineers and clinicians.
\end{itemize}

\cventry{Engineering Internship, Vodafone (Telecom) Headquarters, Accra, Ghana}{}
\begin{itemize}
  \item Designed and proposed a telemedicine smartphone application concept for patient monitoring.
  \item Assisted in vendor management and quality audits for customer-facing telecom hardware systems.
\end{itemize}

\section*{Professional Experience}

\cventry{Senior Manufacturing Engineer, Abbott Laboratories}{2024--2025}
\begin{itemize}
  \item Optimized medical device grinding precision by 31.7\% through DOE, ANOVA, and SPC-driven process improvements.
  \item Applied Monte Carlo simulations, PCA, and variance component analysis to identify and mitigate process variability.
\end{itemize}

\cventry{Quality Engineer, Vastek Inc., San Diego, CA}{2023--2024}
\begin{itemize}
  \item Designed devices and implemented engineering design practices targeting efficiency improvements; worked
  with a cross-functional team of medical professionals and engineers.
  \item Conducted FMEA, Bayesian risk analysis, and fault-tree analysis to quantify cascading failure probabilities.
\end{itemize}

\section*{Volunteering, Community Service \& Projects}
\begin{itemize}
  \item \textbf{Group Homes Volunteer (Autism Support):} assisted residents with learning and recreational activities.
  \item \textbf{MasterCard Foundation Community Projects:} led sanitation, clean-water, and youth education projects.
  \item \textbf{Oforikrom Outreach Project:} organized public school clean-ups, donations, and science tutoring.
  \item \textbf{Vocational Training:} trained four youths in electrical wiring and biofill installation.
  \item \textbf{Personal Tutoring Initiative:} mentored underprivileged students in mathematics; three students gained
  high-school admission and two gained college admission, all first-generation in their families.
  \item Organized community cleaning (gutters, weeds, trash) in Ashaiman-Zenu to reduce contamination and malaria risk.
\end{itemize}

\section*{Awards \& Fellowships}
\begin{itemize}
  \item MasterCard Foundation Scholarship, Arizona State University and KNUST
  \item MORE Fellowship Award, Arizona State University
  \item Vodafone Telecommunications Design Competition, 1st Place
  \item 2nd Place, KNUST Engineering \& Medical School Innovation Challenge
  \item Dean's List, Arizona State University
\end{itemize}

\section*{Leadership \& Mentorship}
\begin{itemize}
  \item Peer Mentor, MasterCard Foundation Scholars (KNUST \& ASU)
  \item Community Project Facilitator, MasterCard Foundation Ghana
  \item Member, Young African Leaders Initiative (YALI)
  \item Club Executive, Biomedical Engineering Society (KNUST \& ASU)
  \item Leader, NSMQ Science \& Math Quiz Team, Kwanyako SHS
\end{itemize}

\section*{Certifications}
\begin{itemize}
  \item Lean Six Sigma (Green Belt)
  \item Design of Experiments (Minitab \& JMP)
  \item Global Leadership \& Project Management (YALI / ASU Global Leadership Academy)
\end{itemize}

\section*{Professional Memberships}
\noindent Society of Manufacturing Engineers (SME) \textbullet\ National Society of Black Engineers (NSBE) \textbullet\
Biomedical Engineering Society (BMES) \textbullet\ ASU Global Leadership Academy \textbullet\ Toastmasters International
\textbullet\ KNUST Peer Counsellors \& MasterCard Peer Mentors \textbullet\ Ghana Engineering Students Association \textbullet\ TraTech

\section*{Mathematical Computing}
\noindent Python (NumPy, SciPy, pandas, statsmodels) \textbullet\ Lean 4 \textbullet\ MATLAB \textbullet\ GNU Octave

\section*{Scientific Software}
\noindent Git/GitHub \textbullet\ GitHub Actions (CI) \textbullet\ pytest \textbullet\ LaTeX (Overleaf) \textbullet\ R

\section*{Scientific Computing}
\noindent Finite-Difference Methods \textbullet\ Monte Carlo Simulation \textbullet\ Euler--Maruyama \& Milstein
Schemes \textbullet\ Particle Methods \textbullet\ Numerical Linear Algebra \textbullet\ Computational Complexity
Analysis \textbullet\ Numerical Verification \textbullet\ Reproducible Research

\section*{Formal Mathematics}
\noindent Lean 4 \textbullet\ formalisation of selected results in stochastic analysis, including Gronwall's
inequality, contraction bounds, and propagation-of-chaos lemmas. Repository includes an automated compilation
check (GitHub Actions, \texttt{lake build}) against current Mathlib releases.

\section*{Engineering Software}
\noindent SolidWorks \textbullet\ ANSYS Fluent \textbullet\ COMSOL Multiphysics \textbullet\ Autodesk Fusion 360
\textbullet\ SimScale \textbullet\ finite element methods

\section*{Technical Skills}

\noindent \textbf{Statistical \& Process Methods:} ANOVA, PCA, DOE, SPC, optimization, GraphPad Prism, Minitab, JMP.\\[3pt]
\textbf{Other:} Fiji/ImageJ, BioRender, Figma, Inkscape, LabVIEW, LTSpice, Windchill, Zotero.

\section*{Languages}
\noindent English \textbullet\ Twi

\end{document}
